\documentclass[10pt,a4paper]{amsart}
\usepackage[margin=19mm]{geometry}
\usepackage{amsmath,amssymb,mathtools}
\usepackage[scr=rsfs,cal=euler]{mathalfa}
\usepackage{xfrac}
\usepackage[unicode,hidelinks,bookmarks=false]{hyperref}
\newcommand{\ie}{i.e.\ }
\newcommand{\cf}{cf.\ }
\newcommand{\resp}{respectively\ }
\newcommand{\Subsection}{\subsection}
\providecommand{\nobreakdash}{\nobreak\hbox{-}}
\setlength{\emergencystretch}{2em}
\multlinegap0pt
\usepackage{xcolor}
\usepackage{tikz}
\newtheorem{thm}{Theorem}
\newcommand{\CC}{{\mathbb C}}
\newcommand{\RR}{{\mathbb R}}
\newcommand{\ddd}{{\overline{d}}}
\newcommand{\ff}{{\overline{f}}}
\newcommand{\gggg}{{\overline{g}}}
\newcommand{\hh}{{\overline{h}}}
\newcommand{\sD}{\mathcal{D}}
\newcommand{\zz}{{\overline{z}}}
\title{Graphical functions in even dimensions}
\author{Michael Borinsky, Oliver Schnetz}
\date{Source: arXiv:2105.05015v3 (CC BY 4.0)}
\begin{document}
\maketitle

\noindent\emph{Source and licence.} Graphical functions in even dimensions. Version arXiv:2105.05015v3 (CC BY 4.0), distributed by arXiv under the Creative Commons Attribution 4.0 International licence, http://creativecommons.org/licenses/by/4.0/, as recorded in the arXiv licence metadata for this exact version and shown on the abstract page https://arxiv.org/abs/2105.05015v3. Full licence notice: \texttt{licenses/arxiv-2105.05015-CC-BY-4.0.txt}.\par\smallskip

\noindent\textit{Editorial scope.} Counted unit: the article's thm thm2a (source0.tex lines 1600--1602). The article's complete original proof of it is delivered with it (source0.tex lines 4742--4791, one \texttt{proof} environment of 50 lines, which the article places away from the statement). No mathematics is written, repaired or re-derived: the statement and the proof are the article's own lines, in the article's own notation, and no retained line is substituted or re-flowed.  Retained uncounted as the source-local context the proof needs: lines 4703--4706.  Nothing was omitted from any retained span, so the package carries no omission record.  No retained line contains a citation, so the wrapper carries no bibliography and no bibliography entry.  No retained line contains a figure environment, an image-inclusion command or drawing code, so no image is delivered and the package declares no asset dependency.  Editorial formatting only, in addition: an AMS \texttt{amsart} preamble that restates the article's own declarations and macros used by the retained lines, namely the environments \texttt{thm} on separate counters, together with the standard packages those lines need.  The retained environments print in this document as Theorem 1 in order of appearance, whereas the article prints the same items with the numbering of its own full document.  The complete native source of the article as distributed in the arXiv e-print for this exact version is bundled unchanged as \texttt{sources/111327/source0.tex}.
\begin{align}\label{eqprop4}
\sD_n\Delta_n&=(\partial_z\partial_\zz)^{n+1}(z-\zz)^n,\nonumber\\
\Delta_nd_n&=\partial_zd_n\partial_\zz.
\end{align}

\begin{thm}\label{thm2a}
The kernel of $\Delta_n$ on $\CC\backslash\RR$ is $d_nh(z)+\ddd_n\hh(\zz)$ for arbitrary (anti-)holomorphic functions $h(z)$ and $\hh(\zz)$.
\end{thm}

\begin{proof}[Proof of Theorem \ref{thm2a}]
From the second identity in (\ref{eqprop4}) we get $\Delta_nd_nh(z)=\partial_zd_n\partial_\zz h(z)=0$. So, $d_nh(z)$ is in the kernel of $\Delta_n$ for any holomorphic function $h(z)$.
By complex conjugation we get $\Delta_n\ddd_n\hh(\zz)=0$.

We need to show that $d_nh(z)+\ddd_n\hh(\zz)$ exhausts the kernel of $\Delta_n$.
To do this we use the first identity in (\ref{eqprop4}) to see that any function $F$ in the kernel of $\Delta_n$ fulfills the identity
$$
(\partial_z\partial_\zz)^{n+1}(z-\zz)^nF=0.
$$
By repeated integration the kernel of $(\partial_z\partial_\zz)^{n+1}$ are polynomials in $z$ of degree $\leq n$ with anti-holomorphic coefficients plus polynomials in $\zz$ of degree $\leq n$
with holomorphic coefficients. Therefore, there are (independent) (anti-)holomorphic functions $h_k(z)$ ($\hh_k(\zz)$) such that $(z-\zz)^nF(z,\zz)=\sum_{k=0}^nh_k(z)\zz^k+\hh_k(\zz)z^k$.
By a linear transformation we equivalently have (anti-)holomorphic functions $f_k(z)$ ($\ff_k(\zz)$) with
$(z-\zz)^nF(z,\zz)=\sum_{k=0}^n(f_{n-k}(z)+\ff_{n-k}(\zz))(z-\zz)^k$. We get
$$
F=\sum_{k=0}^n\frac{f_k(z)+\ff_k(\zz)}{(z-\zz)^k}.
$$
With
$$
h(z)=\frac{(-1)^nn!}{(2n)!}f_n(z)\quad\text{and}\quad\hh(\zz)=\frac{n!}{(2n)!}\ff_n(\zz)
$$
we obtain
$$
F=d_nh(z)+\ddd_n\hh(\zz)+\sum_{k=0}^n\frac{f_k(z)+\ff_k(\zz)-\frac{(n+k)!n!}{(n-k)!k!(2n)!}[(-1)^{n-k}\partial_z^{n-k}f_n(z)+\partial_\zz^{n-k}\ff_n(\zz)]}{(z-\zz)^k}.
$$
The term $k=n$ is zero. For $n=0$ this proves the theorem, so we assume $n\geq1$ in the following.
Introducing new (anti-)homomorphic functions $g_k(z)$ ($\gggg_k(\zz)$) we get that
$$
F=d_nh(z)+\ddd_n\hh(\zz)+G_{n-1}
$$
where
\begin{equation}\label{pf2a1}
G_m=\sum_{k=0}^m\frac{g_k(z)+\gggg_k(\zz)}{(z-\zz)^k}.
\end{equation}
We have $\Delta_nF=\Delta_nG_{n-1}=0$. We show by induction that $\Delta_nG_m=0$ implies $G_m=0$ for all $m=0,\ldots,n-1$.

From $\Delta_nG_m=0$ we get
$$
\sum_{k=0}^m\left(k\frac{g_k'(z)-\gggg_k'(\zz)}{(z-\zz)^{k+1}}+[n(n+1)-k(k+1)]\frac{g_k(z)+\gggg_k(\zz)}{(z-\zz)^{k+2}}\right)=0,
$$
with $n(n+1)-k(k+1)=(n-k)(n+k+1)$. For $m=0$ the above equation implies $G_0=0$. We assume $m>0$ and multiply the above equation with $(z-\zz)^2/(n-m)(n+m+1)$.
In the first term we also shift the summation index $k$ to $k+1$ and obtain
$$
\frac{g_m(z)+\gggg_m(\zz)}{(z-\zz)^m}=-\sum_{k=0}^{m-1}\frac{(k+1)(g_{k+1}'(z)-\gggg_{k+1}'(\zz))+(n-k)(n+k+1)(g_k(z)+\gggg_k(\zz))}{(n-m)(n+m+1)(z-\zz)^k}.
$$
Substitution into (\ref{pf2a1}) gives
$$
G_m=\sum_{k=0}^{m-1}\frac{-(k+1)(g_{k+1}'(z)-\gggg_{k+1}'(\zz))-(m-k)(m+k+1)(g_k(z)+\gggg_k(\zz))}{(n-m)(n+m+1)(z-\zz)^k}.
$$
The right hand side is $G_{m-1}$ for some new functions $\widetilde{g}_k(z)$, $\widetilde{\gggg}_k(\zz)$ and vanishes by induction.
\end{proof}

\end{document}
